\renewcommand{\thetheorem}{17.e.\arabic{theorem}}
\setcounter{theorem}{0}

% ---------------------------------------------------------------------------
% Provenance of the prose in this file.
%   % Extractor: <model>   the block below was transcribed from Prof. Biran's
%                          handwritten notes by that model
%   % Creator:   <model>   the block below was written by that model and has no
%                          counterpart in the notes
% A tag holds until the next tag.
% ---------------------------------------------------------------------------

% Extractor: Gemini 3.1 Pro
\section{Elementary Row Operations Revisited}
\label{sec:eros_revisited}

Let $A \in \M_{n \times p}(K)$. We have previously learned about elementary row operations, which are utilized as an integral part of the Gaussian elimination algorithm. The main point of this section is to demonstrate that all of these operations can be expressed through matrix multiplications.

\begin{notation}
\label{not:matrix_units}
Let $n \in \mathbb{Z}_{\ge 1}$, and let $1 \le i \le n$, $1 \le j \le n$. We denote by $E_{ij} \in \M_{n \times n}(K)$ the matrix whose entries are all zero, except for the entry in the $i$-th row and $j$-th column, which is equal to $1$.
\end{notation}

We will now define three types of elementary matrices in $\M_{n \times n}(K)$. The notes introduce them without a number of their own, so that the numbering of this section starts at \cref{lemma:row_operations_matrix_multiplication} below.

% The notes present the three types under "Notation", unnumbered; keeping this
% definition unnumbered is what makes 17.e.1--17.e.5 below line up with them.
\begin{definition*}[Elementary Matrices]
We define the following three types of elementary matrices:
\begin{enumerate}
    \item \textbf{Type I:} Let $1 \le i \le n$, $1 \le j \le n$ with $i \neq j$, and let $\alpha \in K$. We define $Q_{ij}(\alpha)$ to be the matrix obtained from $I_n$ by applying the row operation $R_i + \alpha R_j \longrightarrow R_i$ (i.e., replacing the $i$-th row by the sum of the $i$-th row and $\alpha$ times the $j$-th row). Explicitly, we have:
    \begin{equation*}
        Q_{ij}(\alpha) = I_n + \alpha \cdot E_{ij}
    \end{equation*}
    
    \item \textbf{Type II:} Let $1 \le i \le n$, $1 \le j \le n$ with $i \neq j$. We define $P_{ij}$ to be the matrix obtained from $I_n$ by permuting (exchanging) the $i$-th row and the $j$-th row, denoted by $R_i \leftrightarrow R_j$.
    
    \item \textbf{Type III:} Let $1 \le i \le n$, and let $0 \neq \alpha \in K$. We denote by $S_i(\alpha)$ the matrix obtained from $I_n$ by applying the row operation $\alpha R_i \longrightarrow R_i$ (i.e., multiplying the $i$-th row by the non-zero scalar $\alpha$).
\end{enumerate}
\end{definition*}

\begin{exercise}[Explicit Form of the Permutation Matrix]
\label{exc:permutation_matrix_form}
Show that the Type II elementary matrix can be written explicitly as $P_{ij} = I_n - E_{ii} - E_{jj} + E_{ij} + E_{ji}$.
\end{exercise}

\begin{example}[Examples of Elementary Matrices]
Consider the space $\M_{4 \times 4}(K)$. Examples of the three types of elementary matrices are:
\begin{align*}
    Q_{1,3}(\alpha) &= \begin{pmatrix} 1 & 0 & \alpha & 0 \\ 0 & 1 & 0 & 0 \\ 0 & 0 & 1 & 0 \\ 0 & 0 & 0 & 1 \end{pmatrix} \\
    P_{2,4} &= \begin{pmatrix} 1 & 0 & 0 & 0 \\ 0 & 0 & 0 & 1 \\ 0 & 0 & 1 & 0 \\ 0 & 1 & 0 & 0 \end{pmatrix} \\
    S_3(\alpha) &= \begin{pmatrix} 1 & 0 & 0 & 0 \\ 0 & 1 & 0 & 0 \\ 0 & 0 & \alpha & 0 \\ 0 & 0 & 0 & 1 \end{pmatrix}
\end{align*}
\end{example}

\begin{lemma}[Matrix Multiplication and Row Operations]
% lemma:17.e.1
\label{lemma:row_operations_matrix_multiplication}
Let $A \in \M_{n \times p}(K)$. Multiplying $A$ from the left by an elementary matrix performs the corresponding elementary row operation on $A$:
\begin{enumerate}
    \item Multiplying $A$ from the left with $Q_{ij}(\alpha)$ results in the row operation $R_i + \alpha R_j \longrightarrow R_i$:
    \begin{equation*}
        A \xrightarrow{R_i + \alpha R_j \longrightarrow R_i} Q_{ij}(\alpha) \cdot A
    \end{equation*}
    \item Multiplying $A$ from the left with $P_{ij}$ results in the row operation $R_i \leftrightarrow R_j$:
    \begin{equation*}
        A \xrightarrow{R_i \leftrightarrow R_j} P_{ij} \cdot A
    \end{equation*}
    \item Multiplying $A$ from the left with $S_i(\alpha)$ results in the row operation $\alpha R_i \longrightarrow R_i$:
    \begin{equation*}
        A \xrightarrow{\alpha R_i \longrightarrow R_i} S_i(\alpha) \cdot A
    \end{equation*}
\end{enumerate}
\end{lemma}

% Creator: Opus 5
% The elementary matrices were defined in \M_{n x n}(K), which is what the
% lemma above needs, since it multiplies A in \M_{n x p}(K) from the left.
% Here they multiply A from the right, so they have to be the p x p ones:
% otherwise none of the three products below is defined, unless n = p.
\begin{importantremark}[Column Operations]
\label{rem:column_operations}
% Extractor: Gemini 3.1 Pro
The same principle holds true for elementary column operations, but via right-multiplication.
% Creator: Opus 5
Here the elementary matrices are the ones of $\M_{p \times p}(K)$, that is, their indices run over $1 \le i \le p$, $1 \le j \le p$: it is the $p$ columns of $A \in \M_{n \times p}(K)$ that are being operated on, and a product such as $A \cdot Q_{ij}(\alpha)$ is defined only for a $p \times p$ factor on the right.
% Extractor: Gemini 3.1 Pro
\begin{enumerate}[label=\textbf{(\alph*')}]
    \item Multiplying $A$ from the right with $Q_{ij}(\alpha)$ results in the column operation $C_j + \alpha C_i \longrightarrow C_j$:
    \begin{equation*}
        A \xrightarrow{C_j + \alpha C_i \longrightarrow C_j} A \cdot Q_{ij}(\alpha)
    \end{equation*}
    \item Multiplying $A$ from the right with $P_{ij}$ results in the column operation $C_i \leftrightarrow C_j$:
    \begin{equation*}
        A \xrightarrow{C_i \leftrightarrow C_j} A \cdot P_{ij}
    \end{equation*}
    \item Multiplying $A$ from the right with $S_i(\alpha)$ results in the column operation $\alpha C_i \longrightarrow C_i$:
    \begin{equation*}
        A \xrightarrow{\alpha C_i \longrightarrow C_i} A \cdot S_i(\alpha)
    \end{equation*}
\end{enumerate}
\end{importantremark}

% Creator: Opus 5
% The single proof below argues (a)--(c) of the lemma and (a')--(c') of the
% remark, with the remark standing between the lemma and it; the title says so,
% since otherwise the reader cannot tell what is being proved.
\begin{proof}[Proof of \cref{lemma:row_operations_matrix_multiplication} and of \cref{rem:column_operations}]
% Extractor: Gemini 3.1 Pro
Statements \textbf{(a)}, \textbf{(b)}, and \textbf{(c)} can be proven directly by calculating the matrix multiplications; the notes leave this calculation to the reader, and it is carried out for \textbf{(a)} among the solutions at the end of the section.

Statements \textbf{(a')}, \textbf{(b')}, and \textbf{(c')} can be deduced from the row operation versions by passing to the transposed matrix. For example, to prove \textbf{(a')}:
\begin{equation*}
    (A \cdot Q_{ij}(\alpha))\transp = Q_{ij}(\alpha)\transp \cdot A\transp = Q_{ji}(\alpha) \cdot A\transp
\end{equation*}
By part \textbf{(a)}, applied to $A\transp \in \M_{p \times n}(K)$ with the roles of $n$ and $p$ interchanged, this is exactly the matrix obtained from $A\transp$ after applying to it the row operation $R_j + \alpha R_i \longrightarrow R_j$.
% Creator: Opus 5
Now the rows of $A\transp$ are the columns of $A$, so performing that row operation on $A\transp$ is the same as performing the column operation $C_j + \alpha C_i \longrightarrow C_j$ on $A$ and transposing the result. Hence the matrix above is equal to:
% Extractor: Gemini 3.1 Pro
\begin{equation*}
    \left( \text{matrix obtained from } A \text{ after applying } C_j + \alpha C_i \longrightarrow C_j \right)\transp
\end{equation*}
Taking the transpose of both sides, it follows that $A \cdot Q_{ij}(\alpha)$ is the matrix obtained from $A$ after the column operation $C_j + \alpha C_i \longrightarrow C_j$. The proofs for \textbf{(b')} and \textbf{(c')} follow a similar logic, using that $P_{ij}\transp = P_{ij}$ and $S_i(\alpha)\transp = S_i(\alpha)$ in place of $Q_{ij}(\alpha)\transp = Q_{ji}(\alpha)$.
\end{proof}

\begin{lemma}[Inverses of Elementary Matrices]
% lemma:17.e.2
\label{lemma:elementary_matrix_inverses}
Every elementary matrix is invertible, and the inverse of each is an elementary matrix of the exact same type:
\begin{equation*}
    Q_{ij}(\alpha)^{-1} = Q_{ij}(-\alpha), \quad P_{ij}^{-1} = P_{ij}, \quad \text{and} \quad S_i(\alpha)^{-1} = S_i(\alpha^{-1})
\end{equation*}
\end{lemma}

% Extractor: Opus 5
\begin{exercise}[Proof of \cref{lemma:elementary_matrix_inverses}]
\label{exc:elementary_matrix_inverses}
The notes leave this proof as an exercise. \textit{(Hint:} Note that $P_{ij} \cdot P_{ij}$ is the matrix obtained from $P_{ij}$ by exchanging rows $i$ and $j$, which brings us immediately back to $I_n$. Apply similar logical steps for the other types.\textit{)}
\end{exercise}

% Creator: Opus 5
% The two theorems below both lean on this, in opposite directions: the proof of
% 17.e.3 opens with "A is invertible, and hence rank(A) = n", and the proof of
% 17.e.4 closes with "A must have rank n, which means that A is invertible".
% The equivalence is nowhere stated in the notes, in either direction: it is a
% step of a transcribed proof being made citable, which is what the ai-family
% is for. It runs on the AI counter, so 17.e.1--17.e.5 are untouched.
\begin{aiproposition}[Invertibility and Full Rank]
\label{aiprop:invertible_iff_full_rank}
Let $A \in \M_{n \times n}(K)$. Then $A$ is invertible if and only if $\rank(A) = n$.
\end{aiproposition}

\begin{proof}
By \cref{lemma:invertible_iff_T_A_iso}, $A$ is invertible if and only if the associated map $T_A: K^n \to K^n$ is an isomorphism. Source and target here have the same dimension $n$, so by \cref{cor:equal_dim_iso} \textbf{(c)} this happens if and only if $T_A$ is surjective, that is, if and only if $\Img(T_A) = K^n$. Since $\Img(T_A)$ is a subspace of $K^n$, \cref{prop:subspace_dimension_constraints} turns that last condition into $\dim \Img(T_A) = n$. Finally, $\dim \Img(T_A) = \rankcol(A) = \rank(A)$ by \cref{rem:rank_notation_and_properties} \textbf{(c)} and \textbf{(a)}.
\end{proof}

% Extractor: Gemini 3.1 Pro
\begin{theorem}[Decomposition into Elementary Matrices]
% thm:17.e.3
\label{thm:elementary_matrix_decomposition}
For every invertible matrix $A \in \GL_n(K)$, there exists an integer $k \ge 1$ and a sequence of elementary matrices $T_1, \dots, T_k$ such that:
\begin{equation*}
    T_k \cdot T_{k-1} \dots T_1 \cdot A = I_n
\end{equation*}
Furthermore, we have $A = T_1^{-1} \dots T_k^{-1}$ and $A^{-1} = T_k \dots T_1$. In particular, any invertible matrix $A$ can be written as a finite product of elementary matrices (recall from \cref{lemma:elementary_matrix_inverses} that each $T_{\ell}^{-1}$ is again an elementary matrix).
\end{theorem}

\begin{proof}
We already know that every matrix $A \in \M_{n \times n}(K)$ can be brought to a row-reduced echelon matrix, which we will call $A'$, using the Gaussian elimination algorithm (\cref{thm:echelon_form_exists}).

In our case, $A$ is invertible, and hence $\rank(A) = n$ by \cref{aiprop:invertible_iff_full_rank}. Because row operations do not change the rank (row-equivalent matrices have the same row space by \cref{prop:row_space_invariance}, hence the same row rank, which is the rank by \cref{rem:rank_notation_and_properties} \textbf{(a)}), we have $\rank(A') = \rank(A)$, which implies that $\rank(A') = n$. The only $n \times n$ row-reduced echelon matrix with rank $n$ is the identity matrix, meaning $A' = I_n$. Indeed, the rank of $A'$ is its number of pivots by \cref{rem:rank_notation_and_properties} \textbf{(b)}, so $A'$ carries $n$ pivots in $n$ columns, one in each; since the pivots move strictly to the right as one descends the rows, the pivot of the $\ell$-th row lies in the $\ell$-th column, and since $A'$ is row-reduced, that pivot is the only non-zero entry of its column.

Since each step in the Gaussian elimination algorithm corresponds to an elementary operation, we can deduce from \cref{lemma:row_operations_matrix_multiplication} that there exist elementary matrices $T_1, \dots, T_k$ (of types I, II, and III) such that:
\begin{equation*}
    T_k \cdot T_{k-1} \dots T_1 \cdot A = I_n
\end{equation*}
% Creator: Opus 5
(If $A = I_n$ to begin with, the algorithm performs no operation at all and this argument produces no matrices; the statement still holds as written, with $k := 1$ and $T_1 := S_1(1) = I_n$, which is an elementary matrix of Type III because $1 \neq 0$.)
% Extractor: Gemini 3.1 Pro
By left-multiplying successively with the inverses of these elementary matrices, which exist by \cref{lemma:elementary_matrix_inverses}, we isolate $A$ and deduce that $A = T_1^{-1} \dots T_k^{-1}$. Taking the inverse of $A$ yields $A^{-1} = T_k \dots T_1$, since inverting a product reverses its order and undoes each factor, by \cref{prop:gl_group_structure}.
\end{proof}

\begin{theorem}[Matrix Inversion Algorithm]
% thm:17.e.4
\label{thm:matrix_inversion_algorithm}
Let $A \in \M_{n \times n}(K)$. Construct the augmented matrix by placing $A$ alongside $I_n$, denoted as $(A \mid I_n) \in \M_{n \times 2n}(K)$, that is, the matrix whose first $n$ columns are those of $A$ and whose last $n$ columns are those of $I_n$. Then, $A$ is invertible if and only if $(A \mid I_n)$ can be brought, via a sequence of elementary row operations, to the form $(I_n \mid B)$ for some matrix $B \in \M_{n \times n}(K)$. Furthermore, if this is the case, we have $A^{-1} = B$.
\end{theorem}

Before proving this theorem, we require an intermediate technical result.

\begin{lemma}[Block Matrix Multiplication]
% lemma:17.e.5
\label{lemma:block_matrix_multiplication}
Let $B, C, D \in \M_{n \times n}(K)$. Then:
\begin{equation*}
    B \cdot \left(C \mid D\right) = \left(B \cdot C \mid B \cdot D\right)
\end{equation*}
\end{lemma}

% Extractor: Opus 5
\begin{exercise}[Proof of \cref{lemma:block_matrix_multiplication}]
\label{exc:block_matrix_multiplication}
The notes leave this proof as an exercise: it follows directly from the definition of matrix multiplication. \textit{(Hint:} Consider the $(k,\ell)$-th entry of $B \cdot (C \mid D)$ and show that it matches the $(k,\ell)$-th entry of $(B \cdot C \mid B \cdot D)$.\textit{)}
\end{exercise}

\begin{proof}[Proof of \Cref{thm:matrix_inversion_algorithm}]
Assume that $A$ is invertible. By \Cref{thm:elementary_matrix_decomposition}, there exist elementary matrices $T_1, \dots, T_k$ such that $T_k \dots T_1 \cdot A = I_n$. Using \cref{lemma:block_matrix_multiplication}, we can multiply the augmented matrix $(A \mid I_n)$ from the left by this chain of elementary matrices:
\begin{equation*}
    T_k \dots T_1 \cdot \left(A \mid I_n\right) = \left(T_k \dots T_1 \cdot A \mid T_k \dots T_1 \cdot I_n\right) = \left(I_n \mid T_k \dots T_1\right)
\end{equation*}
% Creator: Opus 5
Left-multiplying by $T_1$, then by $T_2$, and so on up to $T_k$ is, by \cref{lemma:row_operations_matrix_multiplication}, nothing other than performing the corresponding $k$ elementary row operations in that order. So $(A \mid I_n)$ can indeed be brought to the required form by a sequence of elementary row operations.
% Extractor: Gemini 3.1 Pro

Let us define $B := T_k \dots T_1$. We can clearly see that:
\begin{equation*}
    B \cdot A = (T_k \dots T_1) \cdot A = I_n
\end{equation*}
and using the decomposition of $A$ provided by \cref{thm:elementary_matrix_decomposition},
\begin{equation*}
    A \cdot B = (T_1^{-1} \dots T_k^{-1}) \cdot (T_k \dots T_1) = I_n
\end{equation*}
This explicitly implies that $A^{-1} = B$.

Conversely, assume that $(A \mid I_n)$ can be brought to $(I_n \mid B)$ by a sequence of elementary row operations, for some matrix $B \in \M_{n \times n}(K)$. This directly implies that the matrix $A$ itself is row-equivalent to $I_n$: each of those operations is left-multiplication by an elementary matrix (\cref{lemma:row_operations_matrix_multiplication}), and by \cref{lemma:block_matrix_multiplication} such a multiplication acts on the two blocks separately, so the left-hand block undergoes on its own exactly the same sequence of row operations, ending at $I_n$. Consequently, $A$ must have rank $n$, again because row operations leave the row space, and with it the rank, unchanged (\cref{prop:row_space_invariance} and \cref{rem:rank_notation_and_properties} \textbf{(a)}), which inherently means that $A$ is invertible by \cref{aiprop:invertible_iff_full_rank}. This concludes the proof.
\end{proof}

\begin{example}[Computing the Inverse of a Two-by-Two Matrix]
Let us find the inverse of the matrix $A = \begin{pmatrix} 1 & 2 \\ 3 & 4 \end{pmatrix} \in \M_{2 \times 2}(\mathbb{Q})$. We place $A$ and $I_2$ together in an augmented matrix and apply Gaussian elimination:
\begin{equation*}
    \left( \begin{array}{cc|cc} 1 & 2 & 1 & 0 \\ 3 & 4 & 0 & 1 \end{array} \right) \xrightarrow{R_2 - 3R_1 \longrightarrow R_2} \left( \begin{array}{cc|cc} 1 & 2 & 1 & 0 \\ 0 & -2 & -3 & 1 \end{array} \right)
\end{equation*}
Next, we scale the second row to create a pivot of $1$:
\begin{equation*}
    \xrightarrow{-\frac{1}{2}R_2 \longrightarrow R_2} \left( \begin{array}{cc|cc} 1 & 2 & 1 & 0 \\ 0 & 1 & \frac{3}{2} & -\frac{1}{2} \end{array} \right)
\end{equation*}
Finally, we eliminate the entry above the second pivot:
\begin{equation*}
    \xrightarrow{R_1 - 2R_2 \longrightarrow R_1} \left( \begin{array}{cc|cc} 1 & 0 & -2 & 1 \\ 0 & 1 & \frac{3}{2} & -\frac{1}{2} \end{array} \right)
\end{equation*}
The left side is now the identity matrix $I_2$, which means the right side is our inverse matrix. We conclude that:
\begin{equation*}
    A^{-1} = \begin{pmatrix} -2 & 1 \\ \frac{3}{2} & -\frac{1}{2} \end{pmatrix}
\end{equation*}
\end{example}

% Creator: Opus 5
\subsection{Solutions to the Exercises}
\label{sec:solutions_eros_revisited}

\begin{exercisesolution}[Proof of \cref{exc:permutation_matrix_form}]
Both sides are $n \times n$ matrices, so it suffices to compare entries. Write $M := I_n - E_{ii} - E_{jj} + E_{ij} + E_{ji}$, and recall that the $(k,\ell)$-entry of $E_{pq}$ is $1$ if $(k,\ell) = (p,q)$ and $0$ otherwise, while the $(k,\ell)$-entry of $I_n$ is $\delta_{k\ell}$. Since $i \neq j$, the four matrix units involved sit in four distinct positions, and the entries of $M$ are:
\begin{itemize}
    \item $M_{kk} = 1$ for $k \notin \{i, j\}$, from $I_n$ alone;
    \item $M_{ii} = 1 - 1 = 0$ and $M_{jj} = 1 - 1 = 0$;
    \item $M_{ij} = 0 + 1 = 1$ and $M_{ji} = 0 + 1 = 1$;
    \item $M_{k\ell} = 0$ in every other position.
\end{itemize}
That is exactly the matrix obtained from $I_n$ by exchanging rows $i$ and $j$: row $i$ of $M$ has its single $1$ in column $j$, row $j$ has its single $1$ in column $i$, and every other row is unchanged. Hence $M = P_{ij}$.
\end{exercisesolution}

\begin{exercisesolution}[Proof of \cref{lemma:row_operations_matrix_multiplication} \textbf{(a)}]
Let $A = (a_{k\ell}) \in \M_{n \times p}(K)$. Using $Q_{ij}(\alpha) = I_n + \alpha E_{ij}$ and distributivity,
\[
Q_{ij}(\alpha) \cdot A = (I_n + \alpha E_{ij}) A = A + \alpha \, E_{ij} A.
\]
Now $E_{ij} A$ has $(k, \ell)$-entry $\sum_{q} (E_{ij})_{kq} a_{q\ell}$, and $(E_{ij})_{kq}$ vanishes unless $k = i$ and $q = j$. So $E_{ij}A$ has the $j$-th row of $A$ sitting in its $i$-th row and zeros everywhere else. Adding $\alpha$ times this to $A$ therefore leaves every row of $A$ untouched except the $i$-th, which becomes
\[
(\text{row } i \text{ of } A) + \alpha \cdot (\text{row } j \text{ of } A),
\]
which is precisely the operation $R_i + \alpha R_j \longrightarrow R_i$. The calculations for $P_{ij}$ and $S_i(\alpha)$ are of the same kind.
\end{exercisesolution}

\begin{exercisesolution}[Proof of \cref{lemma:elementary_matrix_inverses}]
For Type II, the hint does it: $P_{ij} \cdot P_{ij}$ is, by \cref{lemma:row_operations_matrix_multiplication} \textbf{(b)}, the matrix obtained from $P_{ij}$ by exchanging its rows $i$ and $j$, and exchanging them a second time restores $I_n$. So $P_{ij}^{-1} = P_{ij}$.

For Type I, apply \cref{lemma:row_operations_matrix_multiplication} \textbf{(a)} twice, or compute directly using $i \neq j$, which gives $E_{ij} E_{ij} = 0$:
\[
Q_{ij}(\alpha) \, Q_{ij}(-\alpha) = (I_n + \alpha E_{ij})(I_n - \alpha E_{ij}) = I_n - \alpha E_{ij} + \alpha E_{ij} - \alpha^2 E_{ij}E_{ij} = I_n,
\]
and the same computation with $\alpha$ and $-\alpha$ interchanged gives $Q_{ij}(-\alpha) Q_{ij}(\alpha) = I_n$. Hence $Q_{ij}(\alpha)^{-1} = Q_{ij}(-\alpha)$.

For Type III, $S_i(\alpha) S_i(\alpha^{-1})$ scales the $i$-th row of $S_i(\alpha^{-1})$ by $\alpha$, turning its $(i,i)$-entry $\alpha^{-1}$ into $1$ and leaving the other rows alone, so the product is $I_n$; symmetrically for the other order. This uses $\alpha \neq 0$, which is part of the definition of $S_i(\alpha)$. In every case the inverse is an elementary matrix of the same type.
\end{exercisesolution}

\begin{exercisesolution}[Proof of \cref{lemma:block_matrix_multiplication}]
Write $(C \mid D) \in \M_{n \times 2n}(K)$ for the matrix whose first $n$ columns are those of $C$ and whose last $n$ columns are those of $D$, and call these columns $u_1, \dots, u_{2n}$. By the column-wise description of a product (\cref{sec:matrix_multiplication}), multiplying on the left by $B$ acts on each column separately:
\[
B \cdot \begin{pmatrix} | & & | \\ u_1 & \dots & u_{2n} \\ | & & | \end{pmatrix} = \begin{pmatrix} | & & | \\ B u_1 & \dots & B u_{2n} \\ | & & | \end{pmatrix}.
\]
The first $n$ columns of $(C \mid D)$ are the columns $c_1, \dots, c_n$ of $C$, and the last $n$ are the columns $d_1, \dots, d_n$ of $D$. So the first $n$ columns of $B \cdot (C \mid D)$ are $Bc_1, \dots, Bc_n$, that is, the columns of $B \cdot C$, and the last $n$ are $Bd_1, \dots, Bd_n$, the columns of $B \cdot D$. Hence $B \cdot (C \mid D) = (B \cdot C \mid B \cdot D)$.
\end{exercisesolution}

\begin{ainote}
All five numbered statements of the notes are present in the transcript, in the
notes' order and with the notes' proofs, but the printed numbers were each one
too large: the three types of elementary matrix are introduced in the notes
under \qt{Notation}, without a number, while the transcript wrapped them in a
numbered \verb|definition|, which consumed 17.e.1 and pushed everything after
it. That definition is now unnumbered, and 17.e.1 through 17.e.5 line up with
the handwritten ones again.

The notes call the matrix units $E_{ij}$; the transcript renamed them
$R_{ij}$, which collides with the $R_i$ that this very section uses for rows ---
$R_{ij}$ and $R_i + \alpha R_j \to R_i$ appear within two lines of each other.
The notes' $E_{ij}$ is restored. Note that the notes switch from the $E_i$ of
the chapter on linear systems to $R_i$ here for rows; that switch is Prof.
Biran's own and is left alone.

The notes mark four things for the reader: the explicit form of $P_{ij}$, the
direct calculation behind \cref{lemma:row_operations_matrix_multiplication},
\cref{lemma:elementary_matrix_inverses}, and
\cref{lemma:block_matrix_multiplication}. All four are now exercises with
solutions above. The proof of \cref{lemma:block_matrix_multiplication} in the
transcript also carried its opening sentence twice, once with an exercise
prompt and once with a hint.

Finally, the transcript's example of a Type II matrix was $P_{3,4}$, correctly
computed but not the notes' example; it is back to Prof. Biran's $P_{2,4}$.

A later pass looking only for gaps in rigour found one product that does not
compose. The elementary matrices are defined in $\M_{n \times n}(K)$, which is
what \cref{lemma:row_operations_matrix_multiplication} needs, since it
multiplies $A \in \M_{n \times p}(K)$ from the left. But
\cref{rem:column_operations} multiplies the same $A$ from the right by the same
matrices, and $A \cdot Q_{ij}(\alpha)$ is defined only for a $p \times p$
factor. The remark now says that its elementary matrices are the ones of
$\M_{p \times p}(K)$; the mathematics of all three statements is otherwise
untouched, and the proof gains the note that it applies \textbf{(a)} to
$A\transp$ with the roles of $n$ and $p$ interchanged.

The chapter also carries the same pattern seen in the chapter on linear maps
and bases: two of its steps rest on facts the book never
states. \Cref{thm:elementary_matrix_decomposition} opens with \qt{$A$ is
invertible, and hence $\rank(A) = n$}, and the proof of
\cref{thm:matrix_inversion_algorithm} closes with the converse implication;
neither direction is stated anywhere in the notes, each being assembled on the
spot from \cref{lemma:invertible_iff_T_A_iso}, \cref{cor:equal_dim_iso} and
\cref{prop:subspace_dimension_constraints}. Since the equivalence is nowhere in
the lecture material, it is stated as
\cref{aiprop:invertible_iff_full_rank} rather than as an unnumbered
Proposition: the ai-family carries the robot marker that tells a reader this is
not Prof. Biran's, and runs on its own counter, so 17.e.1--17.e.5 stay where
they are.

Seven further steps now name the reason they hold: the existence of the
row-reduced echelon form, the invariance of the rank under row operations
(twice), the identification of the only $n \times n$ row-reduced echelon matrix
of rank $n$, the passage from left-multiplication by a chain of elementary
matrices to a sequence of row operations (twice, once in each direction of
\cref{thm:matrix_inversion_algorithm}), and the reversal of order when a
product is inverted. The statement of
\cref{thm:elementary_matrix_decomposition} asks for $k \ge 1$ while its proof
produces no elementary matrix at all when $A = I_n$; the proof now closes that
gap with $T_1 := S_1(1) = I_n$.

Two smaller things. The proof following \cref{rem:column_operations} argues
both that remark and the lemma above it, so its title now says which; and
where it read \qt{But taking the transpose again}, nothing was in fact being
transposed, the step being the observation that the rows of $A\transp$ are the
columns of $A$. The logic was right either way.
\end{ainote}